% ============================================================
% 中青杯全国大学生数学建模竞赛论文模板
% ============================================================

\documentclass[a4paper]{article}
\usepackage[UTF8]{ctex}
\usepackage{geometry}
\usepackage{amsmath,amssymb}
\usepackage{graphicx}
\usepackage{float}
\usepackage{array}
\usepackage{booktabs}
\usepackage{caption}
\usepackage{fancyhdr}
\usepackage{titlesec}
\usepackage{setspace}
\usepackage{indentfirst}
\usepackage{enumitem}

% ============================================================
% 页面设置
% ============================================================
\geometry{a4paper,left=2.5cm,right=2.5cm,top=2.5cm,bottom=2.5cm}

% 1.25倍行距
\renewcommand{\baselinestretch}{1.25}

% 段落首行缩进2字符
\setlength{\parindent}{2em}
\setlength{\parskip}{0pt}

% ============================================================
% 页眉页脚
% ============================================================
\pagestyle{fancy}
\fancyhf{}
\fancyfoot[C]{\thepage}
\renewcommand{\headrulewidth}{0pt}
\renewcommand{\footrulewidth}{0pt}

\fancypagestyle{plain}{
    \fancyhf{}
    \fancyfoot[C]{\thepage}
    \renewcommand{\headrulewidth}{0pt}
}

% ============================================================
% 标题格式
% ============================================================
\titleformat{\section}
    {\centering\zihao{-3}\bfseries}
    {\chinese{section}、}{0pt}{}
\titlespacing*{\section}{0pt}{20pt}{10pt}

\titleformat{\subsection}
    {\normalsize\bfseries}
    {\thesubsection}{1em}{}
\titlespacing*{\subsection}{0pt}{12pt}{6pt}

\titleformat{\subsubsection}
    {\normalsize\bfseries}
    {\thesubsubsection}{1em}{}
\titlespacing*{\subsubsection}{0pt}{10pt}{5pt}

\renewcommand{\thesection}{\chinese{section}}
\renewcommand{\thesubsection}{\arabic{section}.\arabic{subsection}}
\renewcommand{\thesubsubsection}{\arabic{section}.\arabic{subsection}.\arabic{subsubsection}}

% ============================================================
% 图表标题格式
% ============================================================
\captionsetup{
    labelsep=space,
    justification=centering,
    font=normalsize
}

\setlist{leftmargin=2em,itemindent=0pt,parsep=0pt,topsep=0pt,partopsep=0pt}

% ============================================================
% 正文开始
% ============================================================
\begin{document}

% ============================================================
% 第一页：摘要页
% ============================================================
\thispagestyle{plain}
\setcounter{page}{1}

% 论文题目：三号，不加粗
\begin{center}
    {\zihao{3} 基于K-means和BP神经网络的智能养老辅助系统研究}
\end{center}

% “摘要”二字：四号，不加粗
\begin{center}
    \zihao{4} 摘要
\end{center}

随着全球人口老龄化的加速，智能养老服务的需求日益增加。本文建立了K-means+BP神经网络的组合模型以及支持向量机SVM模型，解决了不同区域养老服务需求预测问题。模型评估采用R²和RMSE，结果验证了模型的可靠性。

\vspace{0.5cm}

\noindent\textbf{关键词：} K-means聚类；BP神经网络；支持向量机

\newpage

% ============================================================
% 正文
% ============================================================

\section{问题重述}

\subsection{问题背景}

人口老龄化是全球性问题。根据联合国数据，全球65岁及以上人口正在迅速增加。

\subsection{问题提出}

如何利用人工智能技术，为不同区域的老年人提供个性化养老服务。

\section{模型假设}

\begin{enumerate}
    \item 假设各地区经济水平保持稳定。
    \item 假设养老服务体系不会发生巨大变化。
    \item 假设数据准确可靠。
\end{enumerate}

\section{符号说明}

\begin{table}[H]
    \centering
    \begin{tabular}{|c|c|}
        \hline
        \textbf{符号} & \textbf{含义} \\
        \hline
        $m_j$ & 簇中心 \\
        $C$ & 惩罚因子 \\
        \hline
    \end{tabular}
\end{table}

\section{模型的建立与求解}

\subsection{K-means模型}

K-means是常用聚类算法。距离公式如下：

\[
distance(x_i,m_j) = \|x_i - m_j\| \tag{1}
\]

\subsection{BP神经网络模型}

BP神经网络用于分类预测。模型结果如表1所示：

\begin{table}[H]
    \centering
    \caption{模型结果}
    \begin{tabular}{|c|c|}
        \hline
        \textbf{R²} & \textbf{RMSE} \\
        \hline
        0.96326 & 1.7042 \\
        \hline
    \end{tabular}
\end{table}

\section{结论}

本文建立的模型具有良好的拟合度和预测准确性。

% ============================================================
% 参考文献
% ============================================================
\begin{thebibliography}{99}
    \bibitem{1} 黄韬.基于k-means聚类算法的研究[J].计算机技术与发展,2011.
    \bibitem{2} 李萍.基于MATLAB的BP神经网络预测系统的设计[J].计算机应用与软件,2008.
\end{thebibliography}

\appendix
\section{附录}

附录内容：相关代码和数据表。

\end{document}
